\documentclass[12pt,a4paper]{article}

\usepackage[utf8]{inputenc}
\usepackage[T1]{fontenc}
\usepackage{lmodern}
\usepackage[french,english]{babel}
\usepackage{geometry}
\geometry{margin=2.5cm}
\usepackage{graphicx}
\usepackage{xcolor}
\usepackage{tikz}
\usetikzlibrary{shapes.geometric, arrows.meta, positioning, calc, decorations.pathmorphing, fit, backgrounds}
\usepackage{pgfplots}
\pgfplotsset{compat=1.18}
\usepackage{hyperref}
\hypersetup{colorlinks=true, linkcolor=blue!70!black, urlcolor=blue!60!black}
\usepackage{listings}
\usepackage{booktabs}
\usepackage{tabularx}
\usepackage{multirow}
\usepackage{float}
\usepackage{caption}
\usepackage{amsmath,amssymb}
\usepackage{enumitem}
\usepackage{fancyhdr}
\usepackage{titlesec}

\definecolor{darkbg}{HTML}{0b0f19}
\definecolor{cyan}{HTML}{22d3ee}
\definecolor{green}{HTML}{34d399}
\definecolor{red}{HTML}{f87171}
\definecolor{amber}{HTML}{fbbf24}
\definecolor{purple}{HTML}{a78bfa}
\definecolor{codebg}{HTML}{1a2332}
\definecolor{codegray}{HTML}{94a3b8}

\lstdefinestyle{python}{
    backgroundcolor=\color{codebg},
    basicstyle=\ttfamily\small\color{white},
    keywordstyle=\color{cyan}\bfseries,
    commentstyle=\color{codegray}\itshape,
    stringstyle=\color{green},
    numberstyle=\tiny\color{codegray},
    breaklines=true, frame=single,
    rulecolor=\color{cyan!30},
    numbers=left, numbersep=8pt, tabsize=4,
    showstringspaces=false
}

\pagestyle{fancy}
\fancyhf{}
\fancyhead[L]{\small\textsc{AEGIS} --- Defense \& Deception Subsystem}
\fancyhead[R]{\small ENSA Kenitra --- AI Module 2026}
\fancyfoot[C]{\thepage}
\renewcommand{\headrulewidth}{0.4pt}

\titleformat{\section}{\Large\bfseries\color{cyan!80!black}}{\thesection}{1em}{}
\titleformat{\subsection}{\large\bfseries\color{cyan!60!black}}{\thesubsection}{1em}{}
\titleformat{\subsubsection}{\normalsize\bfseries\color{cyan!40!black}}{\thesubsubsection}{1em}{}

\begin{document}

\begin{titlepage}
\centering
\vspace*{2cm}
\begin{tikzpicture}
    \node[draw=cyan, line width=2pt, rounded corners=8pt,
          minimum width=12cm, minimum height=4cm, fill=darkbg!3] (box) {};
    \node[anchor=center] at (box.center) {
        \begin{minipage}{10cm}
            \centering
            {\Huge\bfseries\color{cyan!80!black} AEGIS}\\[0.5cm]
            {\Large\color{darkbg!70!black} Adaptive Entropy-based Gateway\\for Intrusion Suppression}\\[0.8cm]
            {\large\color{red} Defense \& Deception Subsystem}\\[0.3cm]
            {\large\color{darkbg!50!black} Moving Target Defense --- Technical Documentation}
        \end{minipage}
    };
\end{tikzpicture}
\vspace{2cm}
{\large\textbf{ENSA Kenitra}\\[0.3cm]
{\large AI Module --- Semester 2 --- 2026}}
\vspace{1.5cm}
\begin{tabular}{rl}
\textbf{Document Author:} & El Yazid\\
\textbf{Team Members:} & Adil DILLEKH (ML), Khadija NAFIA (ML)\\
 & El Yazid (Deception), Anas ELKARTOUTI (Docs), Anas MOULAY (VM)\\
\textbf{Date:} & June 2026\\
\end{tabular}
\vfill
{\small\color{darkbg!40!black} Defense and deception subsystem of the AEGIS network intrusion detection system.}
\end{titlepage}

\tableofcontents
\newpage

% ═══════════════════════════════════════════════════════════
\section{Introduction}
\label{sec:intro}

\subsection{Context}

Network Intrusion Detection Systems (NIDS) face an evolving threat landscape where attackers adapt to static defenses. Traditional signature-based systems fail against novel attacks, while ML-based detectors---though more adaptive---still operate on a static attack surface. This creates an arms race: the defender's infrastructure is fixed, and the attacker has unlimited time to probe, fingerprint, and exploit it.

The \textbf{AEGIS} (Adaptive Entropy-based Gateway for Intrusion Suppression) system addresses this limitation by coupling ML-based detection with a \textbf{Moving Target Defense (MTD)} layer. Rather than passively detecting threats, AEGIS actively reshapes its own network topology to confuse, misdirect, and delay attackers.

\subsection{Objectives}

This document covers the defense and deception subsystem, responsible for:

\begin{itemize}[leftmargin=*, itemsep=4pt]
    \item \textbf{Attacker Redirection} --- Diverting malicious connections to honeypot services
    \item \textbf{Phantom Network Simulation} --- Advertising fake subnets via ICMP/DNS responses
    \item \textbf{Port Mutation} --- Rotating exposed service ports on a configurable timer
    \item \textbf{Traffic Shaping} --- Injecting crafted packets via NetfilterQueue to confuse fingerprinting
    \item \textbf{Auto-Blacklisting} --- Blocking repeat offenders at the iptables/nftables level
    \item \textbf{IP Management} --- Dynamic block/unblock with persistence across restarts
\end{itemize}

\subsection{Document Structure}

Section~\ref{sec:architecture} presents the overall system architecture. Sections~\ref{sec:deception}--\ref{sec:ipm} describe each module in detail. Section~\ref{sec:integration} covers integration with the ML detection pipeline. Section~\ref{sec:evaluation} provides evaluation metrics. Section~\ref{sec:conclusion} concludes.

% ═══════════════════════════════════════════════════════════
\section{System Architecture}
\label{sec:architecture}

\subsection{High-Level Overview}

The defense subsystem operates as a layered pipeline that processes network events from detection through response.

\begin{figure}[H]
\centering
\begin{tikzpicture}[
    node distance=1.6cm,
    block/.style={rectangle, draw=cyan, fill=darkbg!5, text width=3cm,
                  minimum height=0.9cm, align=center, rounded corners=4pt,
                  font=\small\sffamily, line width=1pt},
    decision/.style={diamond, draw=amber, fill=amber!5, text width=2.5cm,
                     aspect=2, align=center, font=\small\sffamily, line width=1pt},
    action/.style={rectangle, draw=red!70, fill=red!5, text width=2.8cm,
                   minimum height=0.8cm, align=center, rounded corners=4pt,
                   font=\small\sffamily, line width=1pt},
    arrow/.style={-{Stealth[length=3mm]}, thick, cyan!70!black},
    darrow/.style={-{Stealth[length=3mm]}, thick, amber!70!black},
    aarrow/.style={-{Stealth[length=3mm]}, thick, red!70!black},
    lbl/.style={font=\tiny\sffamily, color=codegray, midway}
]
    % Pipeline nodes
    \node[block] (cap) {Scapy\\Capture};
    \node[block, right=2cm of cap] (ml) {ML\\Detection};
    \node[decision, right=1.8cm of ml] (dec) {Attack\\Detected?};
    \node[action, above=1.5cm of dec] (redir) {Attacker\\Redirection};
    \node[block, right=1.5cm of dec] (pass) {Allow\\Traffic};
    \node[action, below=1.5cm of dec] (block) {Auto-\\Blacklist};
    \node[action, right=1cm of redir] (phantom) {Phantom\\Network};
    \node[action, below=1cm of block] (mutate) {Port\\Mutation};

    % Arrows
    \draw[arrow] (cap) -- node[lbl,above]{raw pkts} (ml);
    \draw[arrow] (ml) -- node[lbl,above]{features} (dec);
    \draw[darrow] (dec) -- node[lbl,right]{yes} (redir);
    \draw[arrow] (dec) -- node[lbl,above]{no} (pass);
    \draw[aarrow] (dec) -- node[lbl,left]{severity} (block);
    \draw[aarrow] (redir) -- (phantom);
    \draw[aarrow] (block) -- (mutate);
\end{tikzpicture}
\caption{High-level defense pipeline: traffic capture, ML classification, and layered response.}
\label{fig:pipeline}
\end{figure}

\subsection{Data Flow}

Network traffic is captured by a Scapy sniffer that monitors the configured interface. Each packet is processed through a feature extraction pipeline that generates flow-level records. These records are classified by the ML ensemble (Random Forest + XGBoost + Isolation Forest) from the detection subsystem. When an attack is detected, the defense subsystem activates the appropriate response layers.

\subsection{Module Map}

\begin{figure}[H]
\centering
\begin{tikzpicture}[
    module/.style={rectangle, draw=cyan, fill=darkbg!5, text width=3.8cm,
                   minimum height=2.8cm, align=left, rounded corners=6pt,
                   font=\small\sffamily, line width=1.2pt, inner sep=8pt},
    link/.style={-{Stealth[length=2mm]}, thin, cyan!50!black, dashed}
]
    % Central coordinator
    \node[module, draw=purple, fill=purple!5, text width=4.5cm,
          minimum height=3.5cm] (active) at (0,0) {
        \textbf{\color{purple}Active Defense Engine}\\[4pt]
        \texttt{active\_defense.py}\\[2pt]
        \footnotesize Orchestrates all defense modules.
        Receives alerts from the ML pipeline
        and triggers appropriate countermeasures.
    };

    % Surrounding modules
    \node[module, above left=2cm and 1cm of active.north] (deception) {
        \textbf{\color{red}Core Deception}\\[4pt]
        \texttt{core\_deception.py}\\[2pt]
        \footnotesize Attacker redirection,
        honeypot activation,
        phantom network simulation.
    };

    \node[module, above right=2cm and 1cm of active.north] (mutator) {
        \textbf{\color{amber}Network Mutator}\\[4pt]
        \texttt{network\_mutator.py}\\[2pt]
        \footnotesize Port rotation via
        iptables/nftables rules,
        randomized port mapping.
    };

    \node[module, below left=2cm and 1cm of active.south] (shaper) {
        \textbf{\color{green}Traffic Shaper}\\[4pt]
        \texttt{traffic\_shaper.py}\\[2pt]
        \footnotesize NetfilterQueue packet
        injection, protocol-level
        confusion and delay.
    };

    \node[module, below right=2cm and 1cm of active.south] (ipm) {
        \textbf{\color{red!80}IP Manager}\\[4pt]
        \texttt{monitor\_interface.py}\\[2pt]
        \footnotesize Block/unblock IPs,
        iptables blackholing,
        dashboard integration.
    };

    % Links
    \draw[link] (active.north) -- (deception.south);
    \draw[link] (active.north) -- (mutator.south);
    \draw[link] (active.south) -- (shaper.north);
    \draw[link] (active.south) -- (ipm.north);
\end{tikzpicture}
\caption{Module architecture: the Active Defense Engine orchestrates four defense modules.}
\label{fig:modules}
\end{figure}

% ═══════════════════════════════════════════════════════════
\section{Core Deception Engine}
\label{sec:deception}

\subsection{Overview}

The Core Deception module (\texttt{core\_deception.py}) implements three deception layers: attacker redirection to honeypots, phantom network advertisement, and IP blacklisting. It maintains an in-memory set of honeypot-interaction flags and provides both programmatic and REST API interfaces.

\subsection{Attacker Redirection}

When the ML pipeline classifies traffic as an attack, the system uses iptables REDIRECT rules to silently divert the attacker's TCP connections to a decoy honeypot service.

\begin{figure}[H]
\centering
\begin{tikzpicture}[
    node distance=1.5cm,
    attacker/.style={rectangle, draw=red, fill=red!8, text width=2.5cm,
                     minimum height=1cm, align=center, rounded corners=4pt,
                     font=\small\sffamily, line width=1pt},
    server/.style={rectangle, draw=cyan, fill=cyan!8, text width=2.5cm,
                   minimum height=1cm, align=center, rounded corners=4pt,
                   font=\small\sffamily, line width=1pt},
    honeypot/.style={rectangle, draw=amber, fill=amber!8, text width=2.5cm,
                     minimum height=1cm, align=center, rounded corners=4pt,
                     font=\small\sffamily, line width=1pt},
    iptables/.style={rectangle, draw=purple, fill=purple!8, text width=2.8cm,
                     minimum height=1.2cm, align=center, rounded corners=4pt,
                     font=\small\sffamily, line width=1.5pt, dashed},
    arrow/.style={-{Stealth[length=3mm]}, thick},
    blocked/.style={-{Stealth[length=3mm]}, thick, red, densely dashed},
    redirected/.style={-{Stealth[length=3mm]}, thick, amber}
]
    \node[attacker] (atk) at (0,0) {\textbf{Attacker}\\192.168.100.10};
    \node[iptables] (fw) at (4,0) {iptables\\REDIRECT\\rule};
    \node[server] (srv) at (8,1) {\textbf{Real Server}\\192.168.100.20};
    \node[honeypot] (hp) at (8,-1) {\textbf{Honeypot}\\Decoy Service};

    \draw[blocked] (atk) -- node[above, font=\tiny, red] {blocked} (srv);
    \draw[redirected] (atk) -- (fw);
    \draw[redirected] (fw) to[bend left=20] node[above, font=\tiny, amber] {REDIRECT} (hp);
\end{tikzpicture}
\caption{Attacker redirection: iptables REDIRECT diverts connections from the real server to the honeypot.}
\label{fig:redirection}
\end{figure}

\subsection{Phantom Network}

The phantom network module responds to ICMP echo requests and ARP queries from attacker-controlled IPs with fabricated network information. This creates the illusion of additional hosts, subnets, and services that do not exist.

\begin{itemize}[leftmargin=*, itemsep=3pt]
    \item Responds to ICMP pings with fake IP addresses from a configurable phantom subnet
    \item Generates synthetic ARP replies to populate the attacker's ARP cache with ghost entries
    \item Logs all phantom interactions for telemetry and dashboard display
    \item Configurable phantom subnet prefix (default: \texttt{10.0.0.0/24})
\end{itemize}

\subsection{Honeypot Interaction Tracking}

The module maintains an in-memory set \texttt{\_honeypot\_hits} that tracks which source IPs have interacted with any honeypot listener. A convenience function \texttt{get\_honeypot\_flag(src\_ip)} returns 1 if the IP has touched a honeypot, 0 otherwise---used as an additional feature for the ML classification pipeline.

\subsection{Logging}

Every deception event is dual-logged:

\begin{enumerate}[leftmargin=*, itemsep=3pt]
    \item \textbf{JSON file}: \texttt{data/deception\_events.jsonl} (one JSON object per line, for programmatic consumption)
    \item \textbf{REST API}: Events are pushed to the Flask dashboard via Socket.IO for real-time display
\end{enumerate}

Each event record contains:
\begin{lstlisting}[style=python, caption={Deception event record schema}]
{
    "timestamp": "2026-06-14T15:30:00",
    "event_type": "REDIRECT | PHANTOM | BLACKLIST",
    "attacker_ip": "192.168.100.10",
    "details": "iptables REDIRECT to port 8080",
    "source_ip": "192.168.100.10"
}
\end{lstlisting}

% ═══════════════════════════════════════════════════════════
\section{Network Mutator}
\label{sec:mutator}

\subsection{Overview}

The Network Mutator (\texttt{network\_mutator.py}) implements the Moving Target Defense principle of \emph{port mutation}---periodically changing which ports are exposed to the network. This defeats port-based fingerprinting and service enumeration.

\subsection{Algorithm}

The mutator operates on a configurable rotation interval (default: 30 seconds). At each rotation:

\begin{enumerate}[leftmargin=*, itemsep=3pt]
    \item Select $k$ ports from a predefined service pool $P = \{p_1, p_2, \ldots, p_n\}$
    \item Generate iptables/nftables rules that forward traffic from old ports to new ports
    \item Apply the rules atomically (remove old, add new) to avoid downtime
    \item Log the rotation event with old $\to$ new port mapping
\end{enumerate}

\subsection{Port Mapping Strategy}

\begin{figure}[H]
\centering
\begin{tikzpicture}[
    node distance=1.2cm,
    phase/.style={rectangle, draw=cyan, fill=darkbg!5, text width=3.5cm,
                  minimum height=1.8cm, align=center, rounded corners=4pt,
                  font=\small\sffamily, line width=1pt},
    arrow/.style={-{Stealth[length=3mm]}, thick, cyan!70!black},
    lbl/.style={font=\tiny\sffamily, color=codegray}
]
    \node[phase] (p1) at (0,0) {
        \textbf{Phase $t$}\\[2pt]
        \footnotesize
        80 $\to$ 8080\\
        443 $\to$ 8443\\
        22 $\to$ 2222
    };
    \node[phase] (p2) at (5,0) {
        \textbf{Phase $t{+}1$}\\[2pt]
        \footnotesize
        80 $\to$ 9080\\
        443 $\to$ 9443\\
        22 $\to$ 2223
    };
    \node[phase] (p3) at (10,0) {
        \textbf{Phase $t{+}2$}\\[2pt]
        \footnotesize
        80 $\to$ 10080\\
        443 $\to$ 10443\\
        22 $\to$ 2224
    };

    \draw[arrow] (p1) -- node[lbl,above]{30s} (p2);
    \draw[arrow] (p2) -- node[lbl,above]{30s} (p3);

    \draw[decorate, decoration={brace, amplitude=6pt, mirror}, thick, red]
        (p1.south west) -- (p3.south east)
        node[midway, below=8pt, font=\small\sffamily, red] {Attacker sees constantly changing ports};
\end{tikzpicture}
\caption{Port rotation across time phases.}
\label{fig:port_rotation}
\end{figure}

\subsection{Implementation Details}

\begin{itemize}[leftmargin=*, itemsep=3pt]
    \item Uses \texttt{iptables -t nat -A PREROUTING} for transparent port forwarding
    \item Supports both \texttt{iptables} and \texttt{nftables} backends
    \item Rotation is performed in a background thread to avoid blocking the main event loop
    \item Port selections are randomized with a seed-based PRNG for reproducibility during testing
    \item Grace period after rotation allows existing connections to drain before rules change
\end{itemize}

% ═══════════════════════════════════════════════════════════
\section{Traffic Shaper}
\label{sec:shaper}

\subsection{Overview}

The Traffic Shaper (\texttt{traffic\_shaper.py}) operates at the packet level using NetfilterQueue (NFQUEUE). It intercepts packets in the kernel's networking stack and applies transformations before forwarding.

\subsection{Packet Processing Pipeline}

\begin{figure}[H]
\centering
\begin{tikzpicture}[
    node distance=1.4cm,
    step/.style={rectangle, draw=cyan, fill=darkbg!5, text width=3cm,
                 minimum height=1cm, align=center, rounded corners=4pt,
                 font=\small\sffamily, line width=1pt},
    kernel/.style={rectangle, draw=green!60, fill=green!5, text width=3cm,
                   minimum height=1.5cm, align=center, rounded corners=6pt,
                   font=\small\sffamily, line width=1.5pt},
    arrow/.style={-{Stealth[length=3mm]}, thick, cyan!70!black}
]
    \node[kernel] (nfq) at (0,0) {NetfilterQueue\\(Kernel)};
    \node[step, right=2cm of nfq] (hook) {Packet\\Hook};
    \node[step, right=1.5cm of hook] (classify) {Classify\\Inbound/Outbound};
    \node[step, right=1.5cm of classify] (mutate) {Mutate\\Payload};
    \node[kernel, right=1.5cm of mutate] (fwd) {Forward\\to Network};

    \draw[arrow] (nfq) -- node[above, font=\tiny] {NFQ} (hook);
    \draw[arrow] (hook) -- (classify);
    \draw[arrow] (classify) -- (mutate);
    \draw[arrow] (mutate) -- (fwd);
\end{tikzpicture}
\caption{NetfilterQueue packet processing pipeline.}
\label{fig:nfq}
\end{figure}

\subsection{Packet Mutation Techniques}

The traffic shaper applies several transformations to confuse attacker fingerprinting:

\begin{table}[H]
\centering
\caption{Packet mutation techniques implemented in the traffic shaper.}
\label{tab:mutations}
\begin{tabularx}{\textwidth}{lX}
\toprule
\textbf{Technique} & \textbf{Description} \\
\midrule
TTL Jitter & Randomizes the IP TTL field by $\pm 1$--$3$ hops, confusing traceroute-based fingerprinting \\
Window Size Noise & Adds small random offsets to the TCP window size field \\
Decoy Injection & Injects fake SYN/ACK packets with randomized source IPs from the phantom subnet \\
Payload Padding & Appends null bytes to short packets to alter their apparent size \\
\bottomrule
\end{tabularx}
\end{table}

\subsection{NFQUEUE Configuration}

The traffic shaper uses the following nftables rule to queue inbound TCP traffic:

\begin{lstlisting}[style=python, caption={NFQUEUE rule creation}]
# Setup NFQUEUE hook
subprocess.run([
    "sudo", "nft", "add", "rule", "inet", "filter",
    "input", "queue", "num", "1"
], check=True)

def process_packet(pkt):
    if pkt.haslayer(TCP):
        pkt[TCP].window += random.randint(-32, 32)
        pkt[IP].ttl = max(1, pkt[IP].ttl + random.randint(-2, 2))
    send(pkt, verbose=0)
\end{lstlisting}

\subsection{Prerequisites}

The traffic shaper requires:
\begin{itemize}[leftmargin=*, itemsep=3pt]
    \item \texttt{root} privileges (NetfilterQueue operates in kernel space)
    \item \texttt{libnetfilter-queue-dev} installed (\texttt{apt install libnetfilter-queue-dev})
    \item \texttt{scapy} for packet crafting and injection
    \item A working \texttt{nftables} or \texttt{iptables} installation
\end{itemize}

% ═══════════════════════════════════════════════════════════
\section{IP Management}
\label{sec:ipm}

\subsection{Overview}

The IP Manager (\texttt{monitor\_interface.py}) provides dynamic IP blocking and unblocking capabilities. It maintains a persistent set of blocked IPs and integrates with both the firewall and the dashboard.

\subsection{Blocking Mechanism}

When a source IP is flagged for blocking (either by the Active Defense Engine or manually via the dashboard), the IP Manager applies a two-layer block:

\begin{enumerate}[leftmargin=*, itemsep=3pt]
    \item \textbf{Blackhole route}: \texttt{ip route add blackhole \{IP\}} -- silently drops all packets to/from the IP at the routing level
    \item \textbf{UFW deny}: \texttt{ufw deny from \{IP\}} -- explicit firewall rule as a secondary barrier
\end{enumerate}

\begin{figure}[H]
\centering
\begin{tikzpicture}[
    node distance=1.6cm,
    input/.style={rectangle, draw=cyan, fill=darkbg!5, text width=3cm,
                  minimum height=0.9cm, align=center, rounded corners=4pt,
                  font=\small\sffamily, line width=1pt},
    block/.style={rectangle, draw=red, fill=red!5, text width=3cm,
                  minimum height=0.9cm, align=center, rounded corners=4pt,
                  font=\small\sffamily, line width=1pt},
    state/.style={circle, draw=amber, fill=amber!5, minimum size=1.2cm,
                  align=center, font=\small\sffamily, line width=1pt},
    arrow/.style={-{Stealth[length=3mm]}, thick, cyan!70!black},
    barrow/.style={-{Stealth[length=3mm]}, thick, red}
]
    \node[input] (detect) {Attack\\Detected};
    \node[state, right=2cm of detect] (check) {IP in\\blocklist?};
    \node[input, above=1.5cm of check] (already) {Skip};
    \node[block, below=1.5cm of check] (block) {Apply\\Block};
    \node[state, right=2cm of block] (persist) {Persist\\to file};
    \node[input, right=2cm of persist] (log) {Log \&\\Dashboard};

    \draw[arrow] (detect) -- (check);
    \draw[arrow] (check) -- node[above, font=\tiny] {yes} (already);
    \draw[barrow] (check) -- node[right, font=\tiny] {no} (block);
    \draw[arrow] (block) -- (persist);
    \draw[arrow] (persist) -- (log);
\end{tikzpicture}
\caption{IP blocking decision flow.}
\label{fig:blocking}
\end{figure}

\subsection{Persistence}

Blocked IPs are persisted to a JSON file (\texttt{data/blocked\_ips.json}) so that blocking survives service restarts. On startup, the IP Manager loads this file and re-applies all iptables/ufw rules.

\subsection{Unblocking}

IPs can be unblocked via:
\begin{itemize}[leftmargin=*, itemsep=3pt]
    \item The Flask dashboard (REST API: \texttt{POST /api/unblock})
    \item Manual command: \texttt{ip route del blackhole \{IP\}} and \texttt{ufw delete deny from \{IP\}}
    \item Automatic expiry (configurable TTL per block entry)
\end{itemize}

\subsection{Dashboard Integration}

The IP Manager exposes its state to the dashboard via Socket.IO events, pushing the current blocklist to all connected clients in real time.

% ═══════════════════════════════════════════════════════════
\section{Active Defense Engine}
\label{sec:active}

\subsection{Overview}

The Active Defense Engine (\texttt{active\_defense.py}) is the central orchestrator that coordinates all defense modules. It receives classified events from the ML pipeline and triggers the appropriate combination of deception, mutation, and blocking actions.

\subsection{Decision Matrix}

The engine applies a severity-based decision matrix:

\begin{table}[H]
\centering
\caption{Defense response decision matrix.}
\label{tab:decision}
\begin{tabularx}{\textwidth}{lcccX}
\toprule
\textbf{Severity} & \textbf{Redirect} & \textbf{Blacklist} & \textbf{Mutate Ports} & \textbf{Description} \\
\midrule
LOW & \checkmark & & & Light redirection to honeypot only \\
MEDIUM & \checkmark & \checkmark & & Redirect + blacklist attacker IP \\
HIGH & \checkmark & \checkmark & \checkmark & Full defense activation \\
\bottomrule
\end{tabularx}
\end{table}

\subsection{Response Flow}

\begin{figure}[H]
\centering
\begin{tikzpicture}[
    node distance=1.4cm,
    event/.style={rectangle, draw=cyan, fill=darkbg!5, text width=2.8cm,
                  minimum height=0.9cm, align=center, rounded corners=4pt,
                  font=\small\sffamily, line width=1pt},
    action/.style={rectangle, draw=green!60, fill=green!5, text width=2.8cm,
                   minimum height=0.9cm, align=center, rounded corners=4pt,
                   font=\small\sffamily, line width=1pt},
    severity/.style={trapezium, draw=amber, fill=amber!5, text width=2.5cm,
                     align=center, trapezium left angle=70, trapezium right angle=110,
                     font=\small\sffamily, line width=1pt},
    arrow/.style={-{Stealth[length=3mm]}, thick, cyan!70!black}
]
    \node[event] (alert) {ML Alert\\(ATTACK)};
    \node[severity, right=1.8cm of alert] (sev) {Severity\\Classification};
    \node[action, above right=1cm and 1.5cm of sev] (redir) {Redirect to\\Honeypot};
    \node[action, right=1.5cm of sev] (bl) {Blacklist\\Source IP};
    \node[action, below right=1cm and 1.5cm of sev] (mut) {Rotate\\Ports};

    \draw[arrow] (alert) -- (sev);
    \draw[arrow, green!60!black] (sev) -- node[above, font=\tiny] {LOW} (redir);
    \draw[arrow, amber] (sev) -- node[above, font=\tiny] {MEDIUM} (bl);
    \draw[arrow, red] (sev) -- node[below, font=\tiny] {HIGH} (mut);
\end{tikzpicture}
\caption{Severity-based response activation.}
\label{fig:response}
\end{figure}

\subsection{Event Bridge Integration}

The defense subsystem communicates with the dashboard through the Event Bridge (\texttt{bridge/event\_bridge.py}), which streams events via Socket.IO:

\begin{itemize}[leftmargin=*, itemsep=3pt]
    \item \texttt{detection\_event}: ML classification results (timestamp, source IP, verdict, confidence)
    \item \texttt{deception\_event}: Deception actions taken (REDIRECT, PHANTOM, BLACKLIST)
    \item \texttt{blocked\_ips\_update}: Current blocklist state for dashboard synchronization
\end{itemize}

% ============================================================
\section{Integration with ML Pipeline}
\label{sec:integration}

\subsection{Data Flow Architecture}

The defense subsystem integrates with the detection pipeline through a well-defined interface. The ML pipeline (\texttt{detection/src/modeltrain.py}) produces classified events that the defense subsystem consumes.

\begin{figure}[H]
\centering
\begin{tikzpicture}[
    node distance=1.5cm,
    ml/.style={rectangle, draw=purple, fill=purple!5, text width=3cm,
               minimum height=1.2cm, align=center, rounded corners=4pt,
               font=\small\sffamily, line width=1pt},
    defn/.style={rectangle, draw=cyan, fill=cyan!5, text width=3cm,
                 minimum height=1.2cm, align=center, rounded corners=4pt,
                 font=\small\sffamily, line width=1pt},
    dash/.style={rectangle, draw=green!60, fill=green!5, text width=3cm,
                 minimum height=1.2cm, align=center, rounded corners=4pt,
                 font=\small\sffamily, line width=1pt},
    arrow/.style={-{Stealth[length=3mm]}, thick}
]
    \node[ml] (model) {ML Models\RF + XGB + IF};
    \node[defn, right=2cm of model] (bridge) {Event Bridge\(Socket.IO)};
    \node[defn, below=1.5cm of bridge] (defense) {Defense\Modules};
    \node[dash, right=2cm of bridge] (dash) {Dashboard\Flask + Vue};

    \draw[arrow, purple] (model) -- node[above, font=\tiny] {predictions} (bridge);
    \draw[arrow, cyan] (bridge) -- node[right, font=\tiny] {alerts} (defense);
    \draw[arrow, green!60!black] (bridge) -- node[above, font=\tiny] {events} (dash);
\end{tikzpicture}
\caption{Integration: ML detection feeds the Event Bridge, which distributes to both defense modules and the dashboard.}
\label{fig:integration}
\end{figure}

\subsection{Classification Output}

The ML pipeline outputs a JSON record per classified flow:

\begin{lstlisting}[style=python, caption={Detection output record}]
{
    "timestamp": "2026-06-14T15:30:00",
    "source_ip": "192.168.100.10",
    "verdict": "ATTACK",
    "model": "xgboost",
    "confidence": 0.987,
    "severity": "HIGH"
}
\end{lstlisting}

\subsection{Severity Mapping}

The Active Defense Engine maps the ML confidence score and model agreement to a severity level:

\begin{table}[H]
\centering
\caption{Severity mapping from ML output to defense response.}
\label{tab:severity}
\begin{tabularx}{\textwidth}{lXX}
\toprule
\textbf{Severity} & \textbf{Condition} & \textbf{Response} \
\midrule
LOW & Single model flags attack, confidence $< 0.85$ & Log + honeypot redirect \
MEDIUM & 2+ models agree, confidence /bin/bash.85hBc-/bin/bash.95$ & Redirect + blacklist \
HIGH & All models agree, confidence $> 0.95$ & Full defense: redirect + blacklist + port mutation \
\bottomrule
\end{tabularx}
\end{table}

\subsection{Honeypot Flag Feature}

The defense subsystem feeds back into the detection pipeline: the \texttt{honeypot\_flag} (0 or 1) is appended to each flow record as an additional feature. If a source IP has previously interacted with a honeypot listener, this flag is set to 1, providing the ML models with an additional signal for classification.

% ============================================================
\section{Evaluation}
\label{sec:evaluation}

\subsection{Experimental Setup}

The defense subsystem was evaluated on a two-machine setup:
\begin{itemize}[leftmargin=*, itemsep=3pt]
    \item \textbf{Ubuntu Server} (192.168.100.20): AEGIS defense + ML models + Dashboard
    \item \textbf{Kali Linux} (192.168.100.10): Attacker machine running Nmap, custom scanners, and attack scripts
    \item \textbf{Network}: VirtualBox Internal Network, 192.168.100.0/24
    \item \textbf{Dataset}: CICIDS2017 (11 features, 9 CSV + 2 MTD placeholders)
\end{itemize}

\subsection{ML Detection Performance}

The detection pipeline achieves the following metrics on the CICIDS2017 test set:

\begin{table}[H]
\centering
\caption{Detection model performance (11 features).}
\label{tab:ml_results}
\begin{tabularx}{\textwidth}{lXXXcc}
\toprule
\textbf{Model} & \textbf{Accuracy} & \textbf{F1-Score} & \textbf{Precision} & \textbf{Recall} & \textbf{FPR} \
\midrule
Random Forest & 99.911\% & 99.920\% & 99.95\% & 99.89\% & 0.04\% \
XGBoost & 99.914\% & 99.923\% & 99.96\% & 99.89\% & 0.03\% \
Isolation Forest & 73.39\% & 24.63\% & 46.50\% & 16.98\% & 24.28\% \
\bottomrule
\end{tabularx}
\end{table}

All models meet the project requirements: F1 $\geq$ 88\%, Accuracy $\geq$ 90\%, FPR $<$ 6\% (Isolation Forest excluded due to expected majority-class anomaly detection behavior).

\subsection{Defense Response Timing}

\begin{figure}[H]
\centering
\begin{tikzpicture}
\begin{axis}[
    ybar,
    bar width=15pt,
    width=0.8\textwidth,
    height=6cm,
    ylabel={Response Time (ms)},
    symbolic x coords={Redirect, Blacklist, Port Rotate, Phantom, Full Stack},
    xtick=data,
    x tick label style={rotate=15, anchor=east, font=\small},
    ymin=0, ymax=600,
    grid=major,
    grid style={dashed, gray!30},
    nodes near coords,
    every node near coord/.append style={font=\tiny, color=gray!70!black},
    legend style={at={(0.5,-0.25)}, anchor=north, font=\small},
    title={Defense Module Response Times},
    title style={font=\small\sffamily}
]
\addplot[fill=cyan!60, draw=cyan] coordinates {
    (Redirect, 45)
    (Blacklist, 120)
    (Port Rotate, 350)
    (Phantom, 80)
    (Full Stack, 520)
};
\end{axis}
\end{tikzpicture}
\caption{Average response times for each defense module (measured on Ubuntu Server).}
\label{fig:timing}
\end{figure}

\subsection{Deception Effectiveness}

We measured the deception system effectiveness against a 5-phase Nmap attack simulation:

\begin{table}[H]
\centering
\caption{Attack phases and defense responses.}
\label{tab:attack_results}
\begin{tabularx}{\textwidth}{lXcc}
\toprule
\textbf{Phase} & \textbf{Attack Type} & \textbf{Detected} & \textbf{Response} \
\midrule
1 & Fast SYN scan (1000 ports) & Yes & REDIRECT + BLACKLIST \
2 & Service version scan & Yes & REDIRECT \
3 & Slow stealth scan (60s) & Yes & BLACKLIST \
4 & OS detection (Nmap -O) & Yes & REDIRECT \
5 & UDP scan & Yes & REDIRECT + BLACKLIST \
\bottomrule
\end{tabularx}
\end{table}

\textbf{Key observations:}
\begin{itemize}[leftmargin=*, itemsep=3pt]
    \item Fast SYN scan triggers the highest severity (HIGH) due to rapid packet rate
    \item Slow stealth scan is detected by the Isolation Forest anomaly detector despite low rate
    \item Honeypot interaction flag increases accuracy of subsequent classifications by 2.3\%
    \item Port rotation prevents the attacker from completing service enumeration across phases
\end{itemize}

\subsection{Comparison: Static vs. MTD Defense}

\begin{figure}[H]
\centering
\begin{tikzpicture}
\begin{axis}[
    ybar,
    bar width=12pt,
    width=0.75\textwidth,
    height=6cm,
    ylabel={Value},
    symbolic x coords={Detection Rate, Attacker Recon Time, False Positive Rate, Service Fingerprinting Success},
    xtick=data,
    x tick label style={rotate=20, anchor=east, font=\small},
    ymin=0, ymax=100,
    grid=major,
    grid style={dashed, gray!30},
    legend style={at={(0.5, 1.02)}, anchor=south, font=\small},
    title={Static Defense vs. MTD Defense},
    title style={font=\small\sffamily}
]
\addplot[fill=red!40, draw=red] coordinates {
    (Detection Rate, 94)
    (Attacker Recon Time, 15)
    (False Positive Rate, 4)
    (Service Fingerprinting Success, 85)
};
\addplot[fill=cyan!60, draw=cyan] coordinates {
    (Detection Rate, 99)
    (Attacker Recon Time, 78)
    (False Positive Rate, 2)
    (Service Fingerprinting Success, 20)
};
\legend{Static, MTD (AEGIS)}
\end{axis}
\end{tikzpicture}
\caption{Comparison: static defense vs. AEGIS with Moving Target Defense enabled.}
\label{fig:comparison}
\end{figure}

% ============================================================
\section{Conclusion and Future Work}
\label{sec:conclusion}

\subsection{Summary}

This document presented the defense and deception subsystem of the AEGIS network intrusion detection system. The key contributions are:

\begin{enumerate}[leftmargin=*, itemsep=4pt]
    \item \textbf{Core Deception Engine} providing attacker redirection, honeypot deployment, and phantom network simulation---creating a multi-layered deception environment that消耗 attacker resources.
    \item \textbf{Network Mutator} implementing Moving Target Defense through periodic port rotation, defeating port-based fingerprinting and service enumeration.
    \item \textbf{Traffic Shaper} operating at the kernel level via NetfilterQueue to inject packet-level confusion that defeats automated scanning tools.
    \item \textbf{Active Defense Engine} orchestrating all modules with severity-based decision logic, ensuring proportional response to different threat levels.
    \item \textbf{IP Manager} providing dynamic blacklisting with persistence and dashboard integration.
\end{enumerate}

\subsection{Results Achieved}

The complete AEGIS system (detection + defense) achieves:
\begin{itemize}[leftmargin=*, itemsep=3pt]
    \item Detection accuracy of 99.914\% (XGBoost) on CICIDS2017 with only 11 features
    \item Average defense response time of 520ms for full-stack activation
    \item 100\% detection rate across all 5 attack phases in the demo scenario
    \item Attacker reconnaissance time increased by 420\% with MTD enabled
    \item Service fingerprinting success reduced from 85\% to 20\%
\end{itemize}

\subsection{Limitations}

\begin{itemize}[leftmargin=*, itemsep=3pt]
    \item The Traffic Shaper requires root privileges and \texttt{libnetfilter-queue}, limiting deployment to Linux environments
    \item Port rotation introduces brief connection drops during rule updates (mitigated by drain period)
    \item Phantom network is limited to ICMP/ARP; sophisticated attackers can detect phantom hosts via TCP probes
    \item The system does not currently implement IP hopping or address space randomization
\end{itemize}

\subsection{Future Work}

\begin{itemize}[leftmargin=*, itemsep=3pt]
    \item \textbf{Adaptive MTD}: Use reinforcement learning to optimize rotation intervals based on threat level
    \item \textbf{Honeypot Fingerprinting}: Deploy realistic service emulators (SSH, HTTP, MySQL) instead of simple TCP listeners
    \item \textbf{Deception Score}: Develop a composite metric quantifying how much attacker resources were wasted
    \item \textbf{Distributed MTD}: Extend port rotation across multiple servers in a cluster
    \item \textbf{IPv6 Support}: Implement address generation for phantom networks using IPv6 sparse addressing
\end{itemize}

% ============================================================
\begin{thebibliography}{10}
\small

\bibitem{nist2020}
NIST, Moving Target Defense,'' Special Publication 800-183, 2020.

\bibitem{jajodia2011}
S. Jajodia, A. K. Ghosh, V. Subrahmanian, and V. Swarup, \emph{Moving Target Defense: Creating Asymmetric Uncertainty for Cyber Threats}. Springer, 2011.

\bibitem{shannon1948}
C. E. Shannon, A Mathematical Theory of Communication,'' \emph{Bell System Technical Journal}, vol.~27, pp.~379--423, 1948.

\bibitem{cicids2017}
I. Sharafaldin, A. H. Lashkari, and A. A. Ghorbani, Toward Generating a New Intrusion Detection Dataset and Intrusion Traffic Characterization,'' in \emph{Proc. ICISSP}, 2018.

\bibitem{scapy2007}
P. Biondi, Scapy: A Powerful Interactive Packet Manipulation Program,'' \emph{Network Security}, 2007.

\bibitem{netfilter}
Netfilter Project, Netfilter Queue Library,'' \url{https://www.netfilter.org/projects/libnetfilter_queue/}.

\end{thebibliography}

\end{document}
